% Team 7 mid-project evaluation, 2 November 2026 (LaTeX/Beamer version of the
% slide deck). Build: cd docs/slides && pdflatex midterm.tex (twice).
\documentclass[aspectratio=169,11pt]{beamer}
\usepackage[T1]{fontenc}
\usepackage{lmodern}
\usepackage{amsmath}
\usepackage{booktabs}
\usepackage{tikz}

\definecolor{navy}{HTML}{14213D}
\definecolor{paper}{HTML}{FBFBF8}
\definecolor{mist}{HTML}{EEF1F4}
\definecolor{accent}{HTML}{C8691F}
\definecolor{blue}{HTML}{2F6FB0}
\definecolor{body}{HTML}{3C4658}
\definecolor{sand}{HTML}{F7EDE3}

\setbeamercolor{background canvas}{bg=paper}
\setbeamercolor{normal text}{fg=body}
\setbeamercolor{frametitle}{fg=navy}
\setbeamercolor{structure}{fg=blue}
\setbeamercolor{block title}{fg=blue,bg=mist}
\setbeamercolor{block body}{fg=body,bg=mist}
\setbeamercolor{block title alerted}{fg=accent,bg=sand}
\setbeamercolor{block body alerted}{fg=body,bg=sand}
\setbeamertemplate{navigation symbols}{}
\setbeamertemplate{blocks}[rounded]
\setbeamerfont{frametitle}{size=\Large,series=\bfseries}
\setbeamertemplate{footline}{%
  \hfill\usebeamercolor[fg]{normal text}\scriptsize\insertframenumber\,/\,\inserttotalframenumber\hspace*{1.2em}\vspace*{0.8em}}
\newcommand{\eyebrow}[1]{{\color{accent}\footnotesize\bfseries\MakeUppercase{#1}}\par\smallskip}
\newcommand{\bignum}[2]{{\color{#1}\Huge\bfseries #2}}

\title{Streaming a live lecture over two weak internet connections}
\subtitle{Send whole video frames, not scattered packets, over two paths at once}
\author{Team 7 \textperiodcentered{} Advanced Computer Networks}
\date{Mid-project evaluation \textperiodcentered{} 2 November 2026}

\begin{document}

{\setbeamercolor{background canvas}{bg=navy}
\setbeamercolor{title}{fg=paper}\setbeamercolor{subtitle}{fg=mist}
\setbeamercolor{author}{fg=mist}\setbeamercolor{date}{fg=mist}
\begin{frame}[plain]
  \vfill
  {\color{accent}\footnotesize\bfseries ADVANCED COMPUTER NETWORKS \textperiodcentered{} TEAM 7 \textperiodcentered{} MID-PROJECT EVALUATION}\par\bigskip
  {\color{paper}\huge\bfseries Streaming a live lecture over two weak internet connections\par}\bigskip
  {\color{mist}\large Send whole video frames, not scattered packets, over two paths at once\par}
  \vfill
  {\color{mist}\footnotesize Multipath QUIC \textperiodcentered{} emulated cellular links \textperiodcentered{} 2 November 2026}
\end{frame}}

\begin{frame}{A teacher's live video needs a steady upload that one phone often cannot give}
  \eyebrow{1 \textperiodcentered{} The problem}
  \begin{columns}[T]
    \begin{column}{0.31\textwidth}
      \begin{block}{\bignum{blue}{2.6 Mbit/s}}
        steady upload needed for a 720p class video (Zoom's figure)
      \end{block}
    \end{column}
    \begin{column}{0.31\textwidth}
      \begin{block}{\bignum{accent}{31--44\%}}
        of mobile tests in India miss a consistent-quality bar, though average uploads (5.9--8.4 Mbit/s) look fine
      \end{block}
    \end{column}
    \begin{column}{0.31\textwidth}
      \begin{block}{\bignum{navy}{2 links}}
        are often available: WiFi + mobile data, or two phones
      \end{block}
    \end{column}
  \end{columns}
  \vfill
  {\scriptsize Sources: Zoom bandwidth requirements; Opensignal India report, Feb 2026 (56--69\% of tests on the three private operators met its bar).}
\end{frame}

\begin{frame}{Multipath protocols split the video packet by packet}
  \eyebrow{2 \textperiodcentered{} What exists today}
  \begin{columns}[T]
    \begin{column}{0.56\textwidth}
      Multipath QUIC can already use two links in one connection. A \textbf{scheduler} picks a link for every packet:
      \begin{itemize}
        \item MinRTT: the faster link first
        \item round-robin: turn by turn
        \item redundant: both links
        \item ECF: by expected completion (Lim et al., 2017)
      \end{itemize}
      \smallskip
      Commercial bonding (LiveU) is closed; research that sends whole frames (MoMQ, 2026) works at a relay, not at the teacher's phone.
    \end{column}
    \begin{column}{0.4\textwidth}
      \begin{alertblock}{The weak spot}
        Packets of one frame land on both links. A frame can be shown only when its \textbf{last} packet arrives, so the slow link decides when every frame appears.
      \end{alertblock}
    \end{column}
  \end{columns}
\end{frame}

\begin{frame}{Treat each video frame as one unit and place it where it finishes first}
  \eyebrow{3 \textperiodcentered{} Our idea}
  \begin{columns}[T]
    \begin{column}{0.31\textwidth}
      \begin{block}{\bignum{blue}{1}\;\; Predict}
        For each new frame, estimate when it would finish on each link.
      \end{block}
    \end{column}
    \begin{column}{0.31\textwidth}
      \begin{block}{\bignum{blue}{2}\;\; Place or split}
        Send the whole frame on the earlier link; split a large frame across both only if that finishes sooner.
      \end{block}
    \end{column}
    \begin{column}{0.31\textwidth}
      \begin{block}{\bignum{blue}{3}\;\; Rescue}
        If a link stalls (no progress for one timeout), resend its unfinished part on the other link.
      \end{block}
    \end{column}
  \end{columns}
  \vfill
  Built in the sender on top of multipath QUIC; each frame part is its own QUIC stream. The parts are known ideas (ECF, DEMS, stream-aware MPQUIC); the new part is doing it at the teacher's uplink and measuring it with deadlines and decoding rules.
\end{frame}

\begin{frame}{A frame's finish time is predictable, so we can balance it}
  \eyebrow{4 \textperiodcentered{} Why we think it works}
  Finish time of a frame part of $S$ bytes on link $p$:
  \[
    T_p \;=\; \underbrace{\frac{U_p + S}{R_p}}_{\text{drain the queue}}
    \;+\; \underbrace{(\mathrm{sRTT}_p - \mathrm{minRTT}_p)}_{\text{wait in the buffer}}
    \;+\; \underbrace{\frac{\mathrm{minRTT}_p}{2}}_{\text{fly one way}}
  \]
  {\small $U_p$: bytes already queued on $p$; $R_p$: measured delivery rate. Write $d_p$ for the two delay terms.}

  \begin{columns}[T]
    \begin{column}{0.56\textwidth}
      \textbf{Split so both parts finish together:} $T_1(S_1) = T_2(S - S_1)$ gives
      \[
        S_1 = \frac{\dfrac{U_2 + S}{R_2} + d_2 - \dfrac{U_1}{R_1} - d_1}{\dfrac{1}{R_1} + \dfrac{1}{R_2}},
        \qquad 0 \le S_1 \le S .
      \]
    \end{column}
    \begin{column}{0.4\textwidth}
      \begin{block}{Worked example (formula only)}
        \small $R_1 = R_2 = 200$ KB/s, $d_1 = 15$ ms, $d_2 = 75$ ms, empty queues, $S = 85$ KB:\\
        $S_1 = 0.485\,\text{s} / 0.01\,\text{s/KB} = 48.5$ KB.\\
        Both parts done at \textbf{258 ms}; one link alone: 440 ms.
      \end{block}
    \end{column}
  \end{columns}
\end{frame}

\begin{frame}{Same video, same links, every scheduler: measured, not assumed}
  \eyebrow{5 \textperiodcentered{} How we test it}
  \small
  \begin{columns}[T]
    \begin{column}{0.31\textwidth}
      \begin{block}{Testbed}
        Two emulated links on one laptop: delay, speed, 200 ms buffer and outages set exactly; later, real two-operator recordings replayed.
      \end{block}
    \end{column}
    \begin{column}{0.31\textwidth}
      \begin{block}{Video}
        A real 60-s NPTEL blackboard lecture, encoded like a live stream: 2.6 Mbit/s, 25 frames/s.
      \end{block}
    \end{column}
    \begin{column}{0.31\textwidth}
      \begin{block}{Score}
        Frame $i$ is shown iff $L_i \le D = 250$ ms, where $L_i = \max_{j} c_j - t_i$ over $i$ and the frames it depends on. Freeze: a gap $\ge 190$ ms.
      \end{block}
    \end{column}
  \end{columns}
  \medskip
  \begin{alertblock}{Five predictions, written down before any run}
    One link is often not enough (H1) \textperiodcentered{} today's schedulers waste capacity when delays differ (H2) \textperiodcentered{} our scheduler freezes less (H3) \textperiodcentered{} it needs no special protocol (H4) \textperiodcentered{} it helps most with big keyframes (H5).\\[2pt]
    {\footnotesize 8 paired runs per comparison: Wilcoxon's smallest $p = 2/2^8 = 0.0078 <$ Holm's first threshold $0.05/6 = 0.0083$.}
  \end{alertblock}
\end{frame}

\begin{frame}{The measurement setup is built and verified}
  \eyebrow{6 \textperiodcentered{} Built and checked so far}
  \begin{columns}[T]
    \begin{column}{0.48\textwidth}
      \begin{block}{Emulator is accurate}
        Delay within 0.4 ms, speed within 0.2\%, buffer $200 \pm 1$ ms. Checking caught two emulator bugs, now fixed.
      \end{block}
      \begin{block}{Protocol fixes}
        Four small changes to the QUIC library (TQUIC), each with a test; one fixes a connection stall in the library itself.
      \end{block}
    \end{column}
    \begin{column}{0.48\textwidth}
      \begin{block}{Multipath works}
        20 of 20 test transfers used both links (each 49--51\% of the data).
      \end{block}
      \begin{block}{Pipeline}
        Video streamer, one-command runs, scoring and statistics code, all with tests.
      \end{block}
    \end{column}
  \end{columns}
\end{frame}

\begin{frame}{Plan to the final demo on 17 November}
  \eyebrow{8 \textperiodcentered{} Next steps}
  \begin{columns}[T]
    \begin{column}{0.31\textwidth}
      \begin{block}{By 11 Nov}
        Ideal-link reference; real recordings at 3 sites. Decide: is there a problem worth solving?
      \end{block}
    \end{column}
    \begin{column}{0.31\textwidth}
      \begin{block}{12--14 Nov}
        Build the frame scheduler; run the full comparison ($\sim$1,240 runs, overnight).
      \end{block}
    \end{column}
    \begin{column}{0.31\textwidth}
      \begin{block}{15--17 Nov}
        Statistics, plots, report; final demo on 17 Nov.
      \end{block}
    \end{column}
  \end{columns}
  \medskip
  \begin{alertblock}{Main risks}
    Real recordings need two phones on two operators and a cloud server (tools ready and tested). The schedule is tight: about 25 hours of runs in four nights.
  \end{alertblock}
\end{frame}

{\setbeamercolor{background canvas}{bg=navy}
\begin{frame}[plain]
  \vfill
  {\color{paper}\Huge\bfseries Thank you. Questions?\par}\bigskip
  {\color{mist}\large Problem: one weak uplink. Existing fix: split packets over two links.\\ Our bet: split by whole frames, timed to finish together.\par}
  \vfill
  {\color{mist}\scriptsize Video: Prof.\ V.\ Balakrishnan, IIT Madras, NPTEL (CC BY-SA). Stack: TQUIC (Tencent), Linux netem/htb, ffmpeg/x264.}
\end{frame}}

\end{document}
